\section{Distinguished monic coefficient rings over unramified DVRs}
\label{sec:unramified-distinguished-monic}\label{sec:singular-distinguished-monic}

The order of the obstruction is determined by the first coefficient
of valuation one.  We prove the result for every characteristic-zero
DVR with uniformizer \(p\), including arbitrary residue fields and
noncomplete bases.  The source lattices and their exact relations are
finite polynomial constructions over the original DVR.  Completion
is used only for the additive detectors.  We include the supported
length construction and every categorical calculation used below.
Throughout, lattice exact sequences split over the coefficient ring;
they need not split equivariantly.

\begin{theorem}\label{thm:unramified-distinguished-monic}\label{thm:singular-distinguished-monic}
Let \(O\) be a characteristic-zero DVR with residue field \(k\) of
characteristic \(p>0\), and suppose that its normalized valuation
satisfies \(v_O(p)=1\).  Let
\[
 F(t)=t^n+\sum_{i=0}^{n-1}a_it^i,
 \qquad n\ge2,\quad a_i\in pO,\quad a_0\in p^2O,
 \qquad A=O[t]/(F).
\]
Set
\[
 r=\min\{1\le i<n:a_i\notin p^2O\},\qquad N=\min(n,pr),
\]
with \(r=\infty\), \(N=n\) if this set is empty.
For \(G=C_p\) and \(G=S_p\), there are homomorphisms
\[
 s_G:\bigoplus_{j\ge1}\mathbb Z/N\longrightarrow K_0(\Lat_A(G)),
 \qquad
 \eta_G:K_0(\Lat_A(G))\longrightarrow\bigoplus_{j\ge1}\mathbb Z/N
\]
with \(\eta_Gs_G=1\), and with the image of \(s_G\) in the kernel
of extension to \(A[1/p]\).  For \(G=S_p\), the map \(\eta_G\)
annihilates every Young permutation lattice.  Thus the distinguished
section classes have exact order \(N\), and \(\Phi_A\) is not
surjective in degree \(p\).
The residue field may be finite, infinite, or imperfect, and \(O\)
need not be complete.  No reducedness of \(A[1/p]\) is assumed.
The section classes admit polynomial presentation matrices over finite
\'etale local extensions of \(O\); when \(k\) is infinite these
extensions may all be taken to be \(O\) itself.
\end{theorem}

\subsection{Length on the supported singularity category}

For a Noetherian ring \(S\), write
\(D_{\mathrm{sg}}(S)=D^b(\operatorname{coh}S)/\operatorname{Perf}(S)\).
All categories below carry their usual differential graded, or equivalently
stable, enhancements.  A superscript \(\mathrm{ic}\) denotes idempotent
completion.

\begin{lemma}\label{lem:singular-monic-supported-length}
Let \(T\) be a complete regular local ring of dimension two, let
\(S=T[t]/(t^N)\), and let \(Z\) be its closed point.  Put
\[
 \mathcal Q_Z=
 \left(D^b_Z(\operatorname{coh}S)/\operatorname{Perf}_Z(S)\right)^{\mathrm{ic}}.
\]
There is a homomorphism
\[
 \lambda_N:K_0(\mathcal Q_Z)\longrightarrow\mathbb Z/N
\]
which sends a finite-length \(S\)-module \(L\) to
\(\operatorname{length}_T(L)\bmod N\).
Moreover \(\mathcal Q_Z\) is equivalent to the idempotent completion
of the full subcategory of \(D_{\mathrm{sg}}(S)\) consisting of objects
perfect outside \(Z\).
\end{lemma}

\begin{proof}
Write
\[
 \mathcal P_S=\operatorname{Perf}_Z(S),\quad
 \mathcal C_S=D^b_Z(\operatorname{coh}S),\quad
 \mathcal P_T=\operatorname{Perf}_Z(T),
\]
and let \(j:\mathcal P_S\to\mathcal C_S\) be the inclusion.
Restriction of scalars gives an exact functor
\(R:\mathcal C_S\to\mathcal P_T\): the restricted complex is bounded
coherent with finite-length cohomology and is perfect because \(T\)
is regular.  There is also an exact functor
\[
 G:\mathcal P_S\longrightarrow\mathcal P_T,
       \qquad G(P)=P\otimes_S^{\mathbf L}T.
\]
This derived reduction is used only on perfect \(S\)-complexes.

For \(0\le a\le N\), define exact functors
\[
 L_a(P)=\operatorname{Res}_T(P\otimes_S^{\mathbf L}t^aS).
\]
They take values in \(\mathcal P_T\): in a bounded finite-projective
model for \(P\), every term becomes finite projective over \(T\),
since \(t^aS\) is \(T\)-free.  Localization verifies the support
condition.  The exact sequences
\(0\to t^{a+1}S\to t^aS\to T\to0\) give natural cofiber sequences
\(L_{a+1}\to L_a\to G\).  Their endpoints are \(L_0=Rj\) and
\(L_N=0\).  Additivity therefore supplies a homotopy of maps of
nonconnective \(K\)-theory spectra
\begin{equation}\label{eq:singular-monic-spectrum-filtration}
 \mathbf K(R)\mathbf K(j)\simeq N\mathbf K(G).
\end{equation}
This assertion is at the level of spectra, rather than only of Euler
characteristics.

Nonconnective localization and Morita invariance give a cofiber sequence
\[
 \mathbf K(\mathcal P_S)\longrightarrow\mathbf K(\mathcal C_S)
      \longrightarrow\mathbf K(\mathcal Q_Z).
\]
Use the homotopy in \eqref{eq:singular-monic-spectrum-filtration} in
the square whose horizontal maps are \(\mathbf K(j)\) and
\(N:\mathbf K(\mathcal P_T)\to\mathbf K(\mathcal P_T)\), and whose
vertical maps are \(\mathbf K(G)\) and \(\mathbf K(R)\).
Its map on cofibers is
\begin{equation}\label{eq:singular-monic-cofiber-map}
 \mathbf K(\mathcal Q_Z)\longrightarrow
 \operatorname{cofib}\bigl(N:\mathbf K(\mathcal P_T)
                                  \to\mathbf K(\mathcal P_T)\bigr).
\end{equation}
The localization and additivity inputs, including the use of the
idempotent-completed quotient, are
\citet[Theorem 9.8, p.~808; the localization consequence on p.~814;
Remark 9.33 and Theorem 9.34, pp.~815--816]{BlumbergGepnerTabuada2013}.

Regularity identifies \(\mathcal P_T\) with
\(D^b_Z(\operatorname{coh}T)\).  Its ordinary cohomological
\(t\)-structure is bounded, with heart the finite-length \(T\)-modules.
Truncation triangles and composition series identify
\(K_0(\mathcal P_T)\) with \(\mathbb Z\) by Euler length, and
\citet[Theorem 1.1, author p.~2]{AntieauGepnerHeller2019} gives
\(K_{-1}(\mathcal P_T)=0\).  The long exact sequence of the target
cofiber in \eqref{eq:singular-monic-cofiber-map} therefore gives
\(\pi_0=\mathbb Z/N\).  The resulting map \(\lambda_N\) has the
stated values because restriction sends a finite-length module to
its underlying \(T\)-module.  This construction is already defined
on \(\mathbf K(\mathcal Q_Z)\); no vanishing of
\(K_{-1}(\operatorname{Perf}_Z(S))\), and no extension from an
uncompleted Grothendieck group, has been assumed.

We verify the last assertion to identify the domain that will occur
in the proof.  Choose regular parameters \(u,v\) of \(T\); they form
an \(S\)-regular sequence because \(S\) is \(T\)-free.  A bounded
coherent complex supported at \(Z\) has a bounded model of
finite-length modules.  Some common powers \(u^a,v^a\) kill all its
terms.  A map to it from a bounded projective complex \(P\) factors
through \(P/(u^a,v^a)P\).  The latter is represented by tensoring
\(P\) with the Koszul complex on \(u^a,v^a\), so it is perfect and
supported at \(Z\).  The supported-derived model and the
Verdier-quotient factorization criterion are
\citet[Lemmas 2.1, 2.5, and 2.6, author pp.~3--5]{Orlov2011}.
They show that the supported quotient is fully faithful in
\(D_{\mathrm{sg}}(S)\).

For density, represent an object perfect outside \(Z\), up to shift,
by a finite maximal Cohen--Macaulay module \(M\).  The hypersurface
description identifies singularity-category maps with stable module
maps \citep[equation (1.1), p.~2, and Appendix B,
pp.~37--40]{BurkeWalker2012}.  The stable endomorphism module of
\(M\) is the quotient of \(\operatorname{Hom}_S(M,M)\) by the image
of \(\operatorname{Hom}_S(M,S)\otimes_SM\).  It is finite, and this
description commutes with localization.  Its support is contained
in \(Z\), so powers \(u^a,v^a\) kill the stable identity.
The cone of multiplication by \(u^a\) on \(M\) is therefore
\(M\oplus M[1]\) in the singularity category.  Multiplication by
\(v^a\) is zero on this cone there, since scalar multiplication is
natural.  Taking its cone produces an object containing \(M\) as
a direct summand and represented by the Koszul complex of \(M\)
on \(u^a,v^a\).  That complex has finite-length cohomology.
Thus the fully faithful functor has dense image, and becomes an
equivalence after idempotent completion.
\end{proof}

\subsection{Parameter choices over the residue field}

Put \(g(X)=(1+X)^p-1\).  We choose finite \'etale local extensions
\(D_j/O\), units \(c_j\in D_j\), and \(q_j\in pD_j\) such that
\begin{equation}\label{eq:singular-monic-constant-root}
 g(q_j)+c_ja_0=0.
\end{equation}
The associated closed points of \(\mathbb A^1_k\) will have
pairwise disjoint \(\mathbb F_p^\times\)-orbits, each orbit being
free.  If \(r<\infty\), put \(\beta=a_r/p\in O^\times\).
In the case \(n=pr\) we also require
\begin{equation}\label{eq:unramified-monic-bad-parameters}
 \bar c_j+(-\bar c_j\bar\beta)^p\ne0.
\end{equation}

Here are the exact constant formulas.  If \(a_0=0\), take
\(q_j=0\) and any chosen unit \(c_j\).  Otherwise write
\(a_0=p^\nu u\), with \(\nu\ge2\), \(u\in O^\times\), and take
\begin{equation}\label{eq:unramified-monic-parameter-formulas}
 q_j=p^{\nu-1}\theta_j,\qquad c_j=-g(q_j)/a_0.
\end{equation}
This quotient is integral.  Except when \(p=2,\nu=2\), its residue
is \(-\bar\theta_j/\bar u\).  In that exceptional case its residue
is \(-\bar\theta_j(\bar\theta_j+1)/\bar u\).

Suppose first that \(k\) is finite.  Choose distinct prime integers
\(m_j>p\), and monic polynomials \(Q_j\in O[Z]\) lifting
irreducible polynomials of degree \(m_j\) over \(k\).  Set
\(D_j=O[\theta_j]/(Q_j)\).  It is finite free, with residue field
\(k_j\) of degree \(m_j\) over \(k\), and is local with maximal
ideal \(pD_j\): every maximal ideal of a finite integral extension
contracts to the maximal ideal of \(O\), and its special fiber is
a field.  Gauss's lemma shows that \(D_j\) is a domain.
It is consequently a DVR.  The residue polynomial is separable,
so \(Q_j'(\theta_j)\) is a unit and \(D_j/O\) is finite \'etale.
These assertions do not require \(O\) to be Henselian.

For \(a_0=0\), choose \(c_j=\theta_j\); otherwise use
\eqref{eq:unramified-monic-parameter-formulas}.
The residue \(\bar c_j\) has degree \(m_j\) over \(k\).
In the exceptional quadratic expression, the degree of
\(k(\bar\theta_j)\) over \(k(\bar\theta_j^2+\bar\theta_j)\)
is at most two and divides the odd prime \(m_j\), hence is one.
In all other cases the full degree is immediate.  Therefore
\eqref{eq:unramified-monic-bad-parameters} holds: its left side is
the value of a nonzero degree-\(p\) polynomial over \(k\), whereas
\(\bar c_j\) has degree greater than \(p\).

These closed points have distinct degrees over \(k\).
If \(a\bar c_j\) is a \(k\)-Galois conjugate of \(\bar c_j\),
with \(a\in\mathbb F_p^\times\), taking the norm from \(k_j\)
to \(k\) gives \(a^{m_j}=1\).  Since
\(\gcd(m_j,p-1)=1\), one has \(a=1\).
This proves the required orbit separation.  The degrees and norms
in this argument are relative to \(k\), not to \(\mathbb F_p\).

Now suppose that \(k\) is infinite.  Take \(D_j=O\).
If \(a_0=0\), the possible residues \(\bar c_j\) are all of
\(k^\times\).  The same is true using
\eqref{eq:unramified-monic-parameter-formulas} outside the case
\(p=2,\nu=2\), because its residual expression is a nonzero linear
function of \(\bar\theta_j\).  In that exceptional case the possible
residues are a nonzero scalar multiple of the Artin--Schreier image
\[
 \{x^2+x:x\in k\}.
\]
This image is infinite: each fiber contains at most two elements,
whereas \(k\) is infinite.  Remove zero and, when necessary, the
finitely many roots excluded by
\eqref{eq:unramified-monic-bad-parameters}.  Inductively choose
countably many remaining elements in distinct
\(\mathbb F_p^\times\)-orbits.  At each step only finitely many
additional elements are excluded.  Every nonzero rational point
has a free scalar orbit.  Lift the selected residues, or their
chosen preimages \(\bar\theta_j\), to \(O\), and use the formulas
above.  This proves all parameter requirements without any
perfection assumption on \(k\).

\subsection{Finite polynomial lattices over the original DVR}\label{subsec:singular-monic-polynomial-sections}

Fix one of the resulting \(D=D_j,c=c_j,q=q_j\), and put
\(v=(1+q)^{p-1}\in D^\times\).  The following finite recursion
is over the original ring \(D\):
\begin{equation}\label{eq:singular-monic-finite-jet}
 \gamma_0=q,\qquad
 \gamma_d=-\frac{[t^d]\bigl(g(\sum_{h<d}\gamma_ht^h)+cF\bigr)}{pv}
 \quad(1\le d<N),\qquad
 \alpha_c(t)=\sum_{d<N}\gamma_dt^d.
\end{equation}
Before degree \(r\), all earlier coefficients lie in \(pD\),
and \(g(pC)\) is divisible by \(p^2\).  Since \(a_d\in p^2O\),
the numerator is in \(p^2D\), giving \(\gamma_d\in pD\).
At degree \(r<\infty\), it gives
\(\bar\gamma_r=-\bar c\bar\beta\ne0\).
Thereafter the residual partial polynomial has \(t\)-order \(r\).
Its \(p\)-th power has no term in degree \(d<N\le pr\), and
\(g\) reduces to that \(p\)-th power.  The degree-\(d\) term
of \(cF\) is also divisible by \(p\), since \(d<n\).
Thus every division by \(p\) in
\eqref{eq:singular-monic-finite-jet} is integral.  If \(r=\infty\),
the first argument applies throughout.

The new coefficient changes degree \(d\) by precisely
\(g'(q)\gamma_d=pv\gamma_d\).  The constant equation is
\eqref{eq:singular-monic-constant-root}.  We obtain an exact
polynomial identity
\begin{equation}\label{eq:singular-monic-polynomial-root}
 g(\alpha_c(t))+cF(t)=t^NU_c(t),\qquad
 \alpha_c,U_c\in D[t],\quad\deg\alpha_c<N.
\end{equation}
Its constant residue is
\begin{equation}\label{eq:singular-monic-unit-residue}
 \bar U_c(0)=
 \begin{cases}
 \bar c,&r=\infty\text{ or }n<pr,\\
 (-\bar c\bar\beta)^p,&pr<n,\\
 \bar c+(-\bar c\bar\beta)^p,&n=pr.
 \end{cases}
\end{equation}
It is nonzero by the parameter choices.  The algebra
\(A_D=D[t]/(F)\) is local with maximal ideal \((p,t)\), because
its special fiber is \((D/p)[t]/t^n\).  Hence \(U_c(t)\) is a
unit in \(A_D\), including when \(D\) is not complete.

Set
\[
 h_c(X,t)=\frac{g(X)-g(\alpha_c(t))}{X-\alpha_c(t)}.
\]
For \(0\le i\le N\), define \(M_i(c)\) by generators \(e,f\)
and the two relations given by the columns of
\begin{equation}\label{eq:singular-monic-factorizations}
 \begin{pmatrix}
 X-\alpha_c(t)&t^i\\-U_c(t)t^{N-i}&h_c(X,t)
 \end{pmatrix}.
\end{equation}
The determinant is \(g+cF\).  Equivalently, the relations are
\begin{equation}\label{eq:singular-monic-relations}
 (X-\alpha_c)e=U_ct^{N-i}f,\qquad t^ie+h_c(X,t)f=0.
\end{equation}
They give an \(A_D\)-basis
\begin{equation}\label{eq:singular-monic-free-basis}
 e,\qquad X^jf\quad(0\le j\le p-2).
\end{equation}
Indeed the first relation specifies \(Xe\), while monicity of
\(h_c\) specifies \(X^{p-1}f\).  These prescriptions define an
operator on the free module with the displayed basis, satisfy both
relations, and give inverse maps to the presented module.  The same
argument first works over \(D[t]\).  The adjugate of the matrix
annihilates its cokernel by \(g+cF\), so after imposing \(F=0\)
the operator \(\sigma=1+X\) defines a \(C_p\)-action.

At \(i=0\) eliminate \(e=-h_cf\); the remaining relation is
\((g+cF)f=0\).  At \(i=N\), use the unit \(U_c\) in \(A_D\)
to eliminate \(f=U_c^{-1}(X-\alpha_c)e\), leaving
\((g+cF)e=0\).  Consequently
\begin{equation}\label{eq:singular-monic-endpoints}
 M_0(c)\simeq M_N(c)\simeq A_D[C_p].
\end{equation}

In these bases define
\[
 \alpha_i:M_i\longrightarrow M_{i+1},\quad
 \alpha_i=\operatorname{diag}(t,I_{p-1}),\qquad
 \beta_i:M_{i+1}\longrightarrow M_i,\quad
 \beta_i=\operatorname{diag}(1,tI_{p-1}).
\]
Substitution in \eqref{eq:singular-monic-relations} proves
equivariance; both composites are multiplication by \(t\).
For \(1\le i<N\), the sequence
\begin{equation}\label{eq:singular-monic-mesh}
 0\longrightarrow M_i
 \xrightarrow{(\beta_{i-1},\alpha_i)}M_{i-1}\oplus M_{i+1}
 \xrightarrow{(\alpha_{i-1},-\beta_i)}M_i\longrightarrow0
\end{equation}
is coefficient-split exact.  Write a coefficient vector as \((a,v)\),
where \(a\) is its first coordinate.  The two maps are
\[
 (a,v)\longmapsto((a,tv),(ta,v)),\qquad
 ((a_L,v_L),(a_R,v_R))\longmapsto
                  (ta_L-a_R,v_L-tv_R).
\]
The kernel equations give exactness.  A coefficient left inverse
of the first map and a coefficient right inverse of the second are
\[
 ((a_L,v_L),(a_R,v_R))\longmapsto(a_L,v_R),\qquad
 (a,v)\longmapsto((0,v),(-a,0)).
\]
These are actual splittings even when \(t\) is a zero-divisor.
With \(x_i(c)=[M_i(c)]-[A_D[C_p]]\), the mesh and endpoints give
\begin{equation}\label{eq:singular-monic-annihilation}
 x_i(c)=i x_1(c),\qquad Nx_1(c)=0.
\end{equation}
Thus the annihilation relation holds over the original coefficient ring.

The characteristic polynomial of the \(X\)-operator is
\[
 (Z-\alpha_c)h_c(Z,t)+U_ct^N=g(Z)+cF(t).
\]
This follows from the companion block for \(h_c\); the coupling
to the first coordinate contributes \(U_ct^N\).  The characteristic
polynomial of \(\sigma\) on \(M_i\) is therefore \(Z^p-1\).
Put \(B=A_D[1/p]\), a commutative Artinian algebra, possibly
nonreduced.  Averaging a \(B\)-linear splitting makes
\(M_i[1/p]\) projective over \(B[C_p]\).  This group algebra is
finite \'etale over \(B\).  Over every residue field of \(B\),
the separable characteristic polynomial \(Z^p-1\), with every
irreducible factor occurring once, gives the regular representation.
On each local Artinian factor of \(B[C_p]\) the projective module
therefore has rank one and is free.  Hence
\begin{equation}\label{eq:singular-monic-generic-zero}
 M_i(c)[1/p]\simeq B[C_p],\qquad x_i(c)[1/p]=0.
\end{equation}

\subsection{The detector after completion}

We verify the categorical hypotheses over the complete coefficient
ring without imposing a restriction on its residue field.
Let \(\widehat O\) be the completion, put
\(\widehat A=\widehat O[t]/F\), and set
\[
 Q=\widehat O[[t,X]],\qquad R=Q/(F,g)=\widehat A[C_p].
\]
The ring \(Q\) is regular Noetherian of dimension three.
Distinguished division gives the displayed quotient, and \(F,g\)
is a regular sequence: \(F\) is nonzero, and the lowest nonzero
\(X\)-coefficient of \(g\) on \(\widehat A[[X]]\) is the
non-zero-divisor \(p\).  Form
\[
 Y=\operatorname{Proj}Q[T_F,T_g]/(FT_F+gT_g),
\]
with the natural maps
\(\beta:\mathbb P^1_R\hookrightarrow Y\) and
\(\pi:\mathbb P^1_R\to\operatorname{Spec}R\).
The generic-hypersurface equivalence and its inverse are
\begin{equation}\label{eq:singular-monic-generic-equivalence}
 \mathcal F=\beta_*\pi^*:D_{\mathrm{sg}}(R)
        \simeq D_{\mathrm{sg}}(Y),\qquad
 \mathcal F^{-1}=R\pi_*L\beta^*(-1)[-1].
\end{equation}
The precise inputs are \citet[Theorem 2.10, p.~8; the paragraph
after Corollary 2.11, p.~9; and the formula preceding (4.3),
p.~19]{BurkeWalker2012}.  The ambient ring is regular Noetherian
of finite dimension and the sequence is regular, as just checked.
No coefficient-field, equicharacteristic, or residue-perfectness
hypothesis is needed.

For every \(\widehat A\)-free cyclic lattice \(E\), its positive
self-Ext groups are cyclic-group cohomology with coefficients
\(\operatorname{Hom}_{\widehat A}(E,E)\).  They are two-periodic,
killed by \(p\), and of bounded finite dimension over \(k\).
The even and odd operator modules are finite over
\(R[\chi_F,\chi_g]\)
\citep[Remark 3.19, p.~17]{BurkeWalker2012}.
Their high-degree tails are therefore finite graded modules over
\[
 (R/pR)[\chi_F,\chi_g],\qquad
 R/pR=k[t,X]/(t^n,X^p),
\]
with bounded Hilbert functions after reducing the nilpotent maximal
ideal of \(R/pR\).  Their projective support is a finite set of
closed points of \(\mathbb P^1_k\).  The stable-Ext comparison
identifies it with the nonperfect locus of the generic-hypersurface
image \citep[Corollaries 3.18 and 3.20, pp.~17--18, and the end of
Section 8.1, pp.~33--34]{BurkeWalker2012}.  This argument works over
an arbitrary field \(k\), including an imperfect one.

Complete the chart \(z=T_F/T_g\) after extension to
\(\widehat D_j=D_j\otimes_O\widehat O\), at the selected residue
point \(z=\bar c_j\).  Write \(z=c_j+w\).
The formal constant root is obtained from
\[
 G(B)=g(pB)/p^2,\qquad
 G(B)+z(a_0/p^2)=0,
 \qquad G'(B)=(1+pB)^{p-1}.
\]
The chosen constant \(q_j/p\) is an exact solution at \(w=0\),
and the derivative is a unit.  Formal implicit-function iteration,
followed by multiplication by \(p\), gives the chosen branch
\(\gamma_0(z)\in p\widehat D_j[[w]]\)
for the root of \(g(\gamma_0)+za_0=0\).
Repeat the finite recursion
\eqref{eq:singular-monic-finite-jet} over \(\widehat D_j[[w]]\).
It gives
\begin{equation}\label{eq:singular-monic-root-factorization}
\begin{gathered}
 g(\alpha(z,t))+zF=t^NU(z,t),\qquad
 U\in\widehat D_j[[w,t]]^\times,\\
 \alpha(c_j,t)=\alpha_{c_j}(t).
\end{gathered}
\end{equation}
All coefficient divisibility checks remain valid modulo \(p\),
regardless of perfection of the residue field.

The root equation defines the closed immersion
\begin{equation}\label{eq:singular-monic-root-section}
 \iota_j:\operatorname{Spec}S_j\longrightarrow
 \operatorname{Spec}\widehat D_j[[w,t,X]]/(g+(c_j+w)F),
 \qquad S_j=\widehat D_j[[w,t]]/(t^N).
\end{equation}
It has Tor dimension one.  In the regular ambient power-series ring,
the quotient by \(X-\alpha\) is the domain \(\widehat D_j[[w,t]]\)
and the potential reduces to the nonzero element \(t^NU\).
To check regularity on the hypersurface directly, an equality
\((X-\alpha)b=(g+zF)d\) forces \(d\) to be divisible by
\(X-\alpha\) after reduction modulo that element.  Cancellation
in the ambient domain then gives \(b\in(g+zF)\).
Derived root pullback is consequently bounded and preserves perfect
objects.
The image of every source lattice is perfect off the target
closed point, by the finite-support argument above.
Lemma~\ref{lem:singular-monic-supported-length} applies with
\(T=\widehat D_j[[w]]\), a complete regular local ring of
dimension two.  That lemma has no residue-field hypothesis and
already includes the idempotent-completion and supported-category
arguments.  Composing its length map with these functors and with
completion defines
\begin{equation}\label{eq:singular-monic-detector}
 \rho_j:K_0(\Lat_A(C_p))\longrightarrow\mathbb Z/N.
\end{equation}
Every step sends arbitrary coefficient-split exact sequences to
triangles.  This is therefore a whole-category additive detector.

There are only finitely many nonzero selected coordinates on any
lattice: its original nonperfect support is finite, and the chosen
closed points are distinct.  Thus \((\rho_j)_j\) takes values in
the direct sum, even when \(k\) is uncountable.
It kills the regular lattice and the trivial lattice.  The former
is perfect; the latter has generic-hypersurface image \(Y/(X)\)
near a selected unit parameter, where \(X\) is regular since the
potential modulo \(X\) is \(zF\ne0\).

\subsection{Evaluation and split sections}

Work over \(\widehat D_j\), abbreviating \(c_j\) by \(c\).
Before imposing \(F\), the source matrix defines a module \(L_i\)
over \(\widehat D_j[[t,X]]\), free over \(\widehat D_j[[t]]\)
in the basis \eqref{eq:singular-monic-free-basis}.  The polynomial
operator extends to the formal-power-series action: modulo \((p,t)\)
it is nilpotent, including at the endpoints by the eliminations
already displayed.  Thus the asserted formal module and basis are
well-defined.  Multiplication by \(F\) on \(L_i\) is injective,
and \(M_i=L_i/FL_i\).

Let \(j_c\) be the closed constant section \(z=c\) of \(Y_{\widehat D_j}\),
with source \(\operatorname{Spec}\widehat D_j[[t,X]]/(g+cF)\).
On this chart the immersion \(\beta\) is defined by \(F\), since
\(g=-zF\).  Its derived pullback of \(j_{c*}L_i\) has no higher
Tor because \(F\) is regular on \(L_i\), and gives \(M_i\)
on the constant section of \(\mathbb P^1_R\).  The twist
\((-1)\) is trivial on that section; its projection to
\(\operatorname{Spec}R\) is the identity and has no higher direct
images.  The inverse in \eqref{eq:singular-monic-generic-equivalence}
therefore gives \(M_i[-1]\), so
\begin{equation}\label{eq:singular-monic-image-shift}
 \mathcal F(M_i)\simeq j_{c*}L_i[1].
\end{equation}

Multiplication by \(X-\alpha_c(t)\) on \(L_i\) has determinant,
up to sign, \(t^NU_c(t)\), by its characteristic polynomial.
It is injective over the domain \(\widehat D_j[[t]]\).
Since \(\alpha(z,t)\) specializes to the exact finite polynomial
\(\alpha_c(t)\), the root pullback of \(j_{c*}L_i\) has no
higher Tor.  It is killed by \(w\) and over
\(S_0=\widehat D_j[[t]]/(t^N)\) has presentation
\begin{equation}\label{eq:singular-monic-root-presentation}
 -U_ct^{N-i}f=0,\qquad t^ie+p v(t)f=0,
 \quad v(t)=(1+\alpha_c(t))^{p-1}\in S_0^\times.
\end{equation}
It is the ideal \(J_i=(p,t^i)\subset S_0\), through
\[
 e\longmapsto p v(t),\qquad f\longmapsto-t^i.
\]
Here is the complete syzygy check.  If \(a p v-bt^i=0\) in
\(S_0\), reduction modulo \(p\) gives
\(b=t^{N-i}u+p b_1\).  The element \(p\) is regular on
\(S_0\), so the original equality gives \(av=b_1t^i\).
Consequently
\[
 (a,b)=(b_1/v)(t^i,pv)+u(0,t^{N-i}).
\]
These are precisely the displayed relation columns, since \(U_c\)
is a unit.  This calculation includes \(i=0,N\).

As an \(S_j\)-module killed by \(w\), the ideal fits in
\[
 0\longrightarrow J_i\longrightarrow S_j/(w)
      \longrightarrow S_j/(p,w,t^i)\longrightarrow0.
\]
The middle term is perfect, because \(w\) is regular on \(S_j\).
The last term has length \(i\) over \(\widehat D_j[[w]]\).
This uses that \(p\) is a uniformizer, and introduces no
ramification multiplier.  The shift \([1]\) in
\eqref{eq:singular-monic-image-shift} reverses the negative sign
of the ideal class.  Hence the detector starting over
\(\widehat D_j\) satisfies
\begin{equation}\label{eq:singular-monic-witness-values}
 \rho_c([M_i])=i\pmod N.
\end{equation}

If \(k\) is infinite, \(D_j=O\), so there is no coefficient-degree
factor.  If \(k\) is finite, the completed extension
\(\widehat D_j/\widehat O\) is Galois with residue extension
\(k_j/k\).  For clarity, this follows directly from Hensel lifting:
the reduction of \(Q_j\) splits into \(m_j\) distinct roots over
\(k_j\); each lifts uniquely in \(\widehat D_j\), giving the
\(m_j\) automorphisms.  Distinct roots have unit differences, and
the Chinese remainder theorem gives
\[
 \widehat D_j\otimes_{\widehat O}\widehat D_j
          \simeq\prod_{\tau\in\operatorname{Gal}(k_j/k)}\widehat D_j.
\]
The polynomial recursion commutes with these automorphisms.
Consequently restriction to \(O\), followed by completion and
extension back to \(\widehat D_j\), splits the witness into
conjugates supported at \(\tau(\bar c_j)\).
Exactly one survives at the selected point.  Thus, with
\(m_j=[D_j:O]\) and \(m_j=1\) in the infinite case, define
\[
 y_j=[\operatorname{Res}_{D_j/O}M_1(c_j)]-m_j[A[C_p]].
\]
The evaluation is
\begin{equation}\label{eq:singular-monic-coordinate-values}
 \rho_i(y_j)=\delta_{ij}\pmod N.
\end{equation}
No residue degree is divided out, so no condition on
\(\gcd(m_j,N)\) is required.  By
\eqref{eq:singular-monic-annihilation}, \(Ny_j=0\), and the
generic regularity argument shows that \(y_j\) is generically zero.
These give the cyclic split section.
For completeness, coefficient restriction does not enlarge the original
support point: the finite \'etale coefficient map has exact pushforward
preserving perfect complexes, flat base change commutes with \(\pi^*\),
and the closed immersion \(\beta_*\) commutes with that base change.
Thus the generic-hypersurface image of a restricted witness is the
finite pushforward of its image before restriction.  Its nonperfect
support maps to its selected closed point over \(k\).  Distinct
degrees in the finite-field case, or distinct rational points in the
infinite-field case, therefore give the stated zero off-diagonal values.

We fix the normalizer convention explicitly.  For
\(a\in\{1,\ldots,p-1\}\), put
\[
 \phi_a(X)=(1+X)^a-1,\qquad
 V_a(X)=\sum_{h=0}^{a-1}(1+X)^{ph}.
\]
Then \(g(\phi_a(X))=V_a(X)g(X)\), with \(V_a(0)=a\).
The formal change
\[
 (X,z)\longmapsto(\phi_a(X),V_a(X)z)
\]
is an automorphism, multiplies the potential by a unit, and commutes
with the generic-hypersurface diagram.  Under its geometric pullback,
support at \(\bar c\) moves to \(a^{-1}\bar c\).

If its induced ring automorphism is \(h_a(\sigma)=\sigma^a\),
geometric pullback of a module is \(R\otimes_{h_a,R}M\).
The identification \(r\otimes m\mapsto h_a^{-1}(r)m\) makes
the generator act by the inverse group twist.  Thus the concrete
twist \(M^{(a)}\) whose generator acts as \(\mathcal T^a\), when
\(\mathcal T\) acts on \(M\), has its support multiplied by \(a\).
These conventions give inverse bijections of
\(\mathbb F_p^\times\), and both preserve the identity label.

The selected nonzero rational points have disjoint free scalar orbits
when \(k\) is infinite.  In the finite case the relative norm
argument above excludes every nonidentity scalar-conjugate collision,
and the distinct degrees exclude collisions between different indices.
Therefore for the concrete twists
\begin{equation}\label{eq:singular-monic-normalizer-values}
 \rho_i(y_j^{(a)})=
 \begin{cases}1,&i=j\text{ and }a=1,\\0,&\text{otherwise}.
 \end{cases}
\end{equation}
In particular the inverse labeling convention in Mackey's formula
does not change the unique surviving term.

Set \(z_j=\operatorname{Ind}_{C_p}^{S_p}y_j\), and
\(\eta_{S_p}=(\rho_i)_i\operatorname{Res}_{C_p}^{S_p}\).
For a double coset outside the normalizer, the intersection of
\(C_p\) with its conjugate is trivial, since it has prime order.
Its Mackey contribution is induced from the underlying coefficient
class of \(y_j\), which is zero by its virtual coefficient rank.
The quotient \(N_{S_p}(C_p)/C_p=\mathbb F_p^\times\) supplies
each scalar twist once.  Therefore
\(\eta_{S_p}(z_j)=e_j\).  Induction preserves the actual
annihilation relation and generic vanishing.  This gives the
symmetric-group section and proves exact order \(N\).

Finally, a degree-\(p\) Young lattice restricts to a trivial cyclic
lattice when its partition is \((p)\); every other Young subgroup
has order prime to \(p\), so its restriction is a sum of regular
cyclic lattices.  Both types are killed.  All sources, meshes and
generic isomorphisms were constructed over the original \(O\).
Only the detectors used completion.  This proves
Theorem~\ref{thm:unramified-distinguished-monic} in full.

\begin{corollary}\label{cor:distinguished-monic-regularity}
Let \(O\) be any characteristic-zero DVR with \(v_O(p)=1\), let \(n\ge2\), and
let \(F=t^n+\sum_{i<n}a_it^i\) with every \(a_i\in pO\).
For \(A=O[t]/(F)\), the following are equivalent:
\[
 \Phi_A\text{ is an isomorphism};\qquad\Phi_A\text{ is surjective};\qquad A\text{ is regular};\qquad
 a_0\in pO\setminus p^2O.
\]
In the regular case \(A\) is a DVR.  In every other case
Theorem~\ref{thm:singular-distinguished-monic} gives the stated
split infinite torsion obstruction already in degree \(p\).
\end{corollary}

\begin{proof}
The algebra \(A\) is finite free over \(O\), has dimension one,
and is local with maximal ideal \((p,t)\), since its special
fiber is \(k[t]/(t^n)\).  If \(a_0=pu\) with \(u\)
a unit, write
\(F=t^n+p(u+\sum_{i\ge1}(a_i/p)t^i)\).  The expression in
parentheses is a unit in \(A\), so \(p\in(t^n)\) and the maximal
ideal is generated by \(t\).  Thus \(A\) is a one-dimensional
regular local ring, hence a DVR; alternatively \(F\) is Eisenstein.
Its polynomial-generation property follows from
Theorem~\ref{thm:regular-local}.
If \(a_0\in p^2O\), then \(F\in(p,t)^2\) in the regular local
ring \(O[t]_{(p,t)}\).  Thus the cotangent space of \(A\) has
dimension two while \(\dim A=1\), so \(A\) is singular.  Here
\(N\ge2\), and the preceding theorem proves nonsurjectivity.
\end{proof}


\begin{corollary}\label{cor:unramified-monic-completed-field}\label{cor:singular-monic-generic-field-order}
Under the hypotheses of Theorem~\ref{thm:unramified-distinguished-monic},
suppose that \(\widehat K[t]/(F)\) is a field, where
\(\widehat K=\operatorname{Frac}(\widehat O)\).
Then \(N=n\), so the split section has countably many generators
of exact order \(n\).  This includes arbitrary residue fields.
\end{corollary}

\begin{proof}
The constant term is nonzero, since \(F\) is irreducible of
degree \(n\ge2\).  Write \(\nu=v_O(a_0)\ge2\).
A complete discretely valued field has a unique extension of its
absolute value to a finite extension; hence that valuation is
Galois invariant \citep[Theorem 4.7 and Corollary 4.8,
pp.~54--55]{WilliamsWuthrich2025}.  The cited statements require
completeness, not finiteness or perfection of the residue field.
Thus every conjugate root of \(F\) has valuation \(\nu/n\).
The elementary symmetric expressions give
\[
 v_O(a_i)\ge\frac{(n-i)\nu}{n}.
\]
If \(r<\infty\), its valuation-one coefficient forces
\(r\ge n(1-1/\nu)\ge n/2\), so \(pr\ge n\).
If \(r=\infty\), \(N=n\) by definition.  This proves the claim.
For a noncomplete base the hypothesis deliberately concerns the
completed generic algebra.
\end{proof}

\begin{example}\label{ex:singular-monic-p2-lower-coefficient}
Let \(p=2\), \(F=t^3+2t+4\), and \(A=O[t]/(F)\), with \(O=\mathbb Z_2\) or \(O=\mathbb Z_{(2)}\).  Then \(r=1\) and \(N=2<n=3\).
Take
\[
 D=O[\theta]/(\theta^3+\theta+1),\quad
 q=2\theta,\quad c=-\theta(\theta+1),\quad s=1+q,
\]
and put
\[
 a(t)=q-\frac{c}{s}t,\qquad
 U(t)=\frac{c^2}{s^2}+ct.
\]
The polynomial defining \(D\) has no root modulo two, hence is
irreducible and separable there.  Thus \(D\) is finite \'etale local
of degree three over \(O\); over \(\mathbb Z_2\) it is the
unramified cubic extension.  Both \(c\) and \(s\) are units.
Direct expansion, using \(q^2+2q+4c=0\), gives
\[
 g(a(t))+cF(t)=t^2U(t),\qquad g(X)=X^2+2X.
\]
Consequently
\[
 \mathcal T=
 \begin{pmatrix}1+a(t)&-t\\ U(t)t&-1-a(t)\end{pmatrix}
 \quad\text{satisfies}\quad
 \mathcal T^2=I-cF(t)I=I
 \quad\text{over }A_D.
\]
Let \(L\) be the corresponding rank-two \(A_D\)-lattice.
The residue \(\bar c=\bar\theta^2+\bar\theta\) has degree three,
so the one-conjugate evaluation above gives
\[
 \left[\operatorname{Res}_{D/O}L\right]-3[A[C_2]]
\]
exact order two and generic value zero.  The same finite formulas
work over \(\mathbb Z_{(2)}\), because \(s\) remains a unit.
Since \(S_2=C_2\), this symmetric-group class also lies outside the
Young span.  No assertion that the generic algebra is a field is needed.
\end{example}

\begin{remark}\label{rem:distinguished-completed-field-hypothesis}
The completed-field condition cannot be replaced by irreducibility
only over \(\Q\) for the localized base.  The polynomial
\(t^3+2t+4\) has none of its possible integer roots
\(\pm1,\pm2,\pm4\), so it is irreducible over \(\Q\).
But \(F(2u)=4(2u^3+u+1)\), and the latter cubic has a simple
root modulo two.  Hensel's lemma gives a root in \(\Q_2\).
Here \(r=1\) and \(N=2<3\).
\end{remark}
